
3.3.
Soient @@M0@@, @@M1@@ deux @@M2@@-anneaux, @@M3@@ un homomorphisme d'anneaux. D'après la définition 3.2.ii), @@M4@@ est un @@M5@@-homomorphisme si et seulement si pour tout @@M6@@, et tout @@M7@@, @@M8@@, ou encore @@M9@@.

3.4. De la relation 3.2.ii), on déduit les relations :
@@M10@@
Elles montrent que l'homomorphisme @@M11@@ définit sur @@M12@@ une nouvelle structure de pré-@@M13@@-anneau ; nous allons voir que ce n'est pas en général une structure de @@M14@@-anneau. Si on désigne par @@M15@@ la substitution sur l'anneau des séries formelles en @@M16@@ à coefficients dans @@M17@@ définie par @@M18@@, on a :
@@M19@@
L'automorphisme @@M20@@ est fonctoriel par rapport à @@M21@@, et définit un automorphisme du foncteur @@M22@@ considéré comme foncteur en groupes abéliens ; à l'aide de cet automorphisme, on peut, par transport de structure, définir sur @@M23@@ une nouvelle structure de foncteur en @@M24@@-anneaux, qui sera telle que pour tout @@M25@@-anneau @@M26@@, l'homomorphisme @@M27@@ soit un @@M28@@-homomorphisme ; nous allons expliciter la multiplication et les opérations @@M29@@ de cette deuxième structure.

3.5. La multiplication est définie par l'homomorphisme de foncteurs :
@@M30@@
où le signe @@M31@@ note la multiplication définie en 2.3. Cette nouvelle multiplication sera notée @@M32@@. Elle est comme d'habitude définie par sa valeur sur @@M33@@~; on a donc :
